\section{Installation or Compilation}
\subsection{Installation}\label{sec:install}

The \gretlpy{} package is available
via the PyPI infrastructure, so installation requires calling:
\begin{verbatim}
python3(.exe) -m pip install gretl4py
\end{verbatim}

In addition, \texttt{wheel}
snapshots are also available at
\url{https://sourceforge.net/projects/gretl/files/gretl4py/snapshots/}.
To install the package, please download the appropriate \texttt{.whl}
file and run the following command:

\begin{verbatim}
python.exe -m pip install gretl4py-0.2-cp313-cp313-win_amd64.whl
\end{verbatim}

Since installing packages via \texttt{pip} can be problematic on some
Linux distributions, we provide a script named \texttt{webinstall.py}
(available at: \url{https://sourceforge.net/p/gretl/gretl4py/ci/master/tree/}),
which downloads a \texttt{wheel} file and installs it in the user-space area.

\begin{figure}[h!tb]
	\centering
	\begin{subfigure}[b]{\textwidth}
		\centering
		\includegraphics[width=\columnwidth,keepaspectratio]{mswindows_install.py}
		\caption{MS~Windows, \python[3.13]}
		\label{fig:install.py-mswindows}
	\end{subfigure}
	
	\begin{subfigure}[b]{\textwidth}
		\centering
		\includegraphics[width=\columnwidth,keepaspectratio]{macos_install.py}
		\caption{macOS, \python[3.12]}
		\label{fig:install.py-macos}
	\end{subfigure}
	
	\caption{Example output from \texttt{install.py}}
	\label{fig:install.py}
\end{figure}

To verify that installation was successful, run the following command:
\begin{center}
\begin{tabular}{rl}
  \hline
  \texttt{python.exe -c "import gretl; gretl.about(True)"} & on MS
                                                             Windows \\
  \texttt{python3 -c "import gretl; gretl.about(True)"} & on macOS or Linux\\
  \hline
\end{tabular}
\end{center}
Doing so should produce output similar to Figure~\ref{fig:mswindows_about}.

\begin{figure}[h!tb]
	\centering
	\includegraphics[scale=0.7]{about}
	\caption{Example output from \texttt{gretl.about(True)}}
	\label{fig:mswindows_about}
\end{figure}


\subsection{Compilation on GNU/Linux}\label{sec:linux-compile}
Before compiling \gretlpy{} on GNU/Linux, ensure that the following
components are installed:

\begin{enumerate}
	\item \pkg{meson}, \pkg{ninja}, \pkg{pkg-config}, and \gretl
		version 2025b-git or later (see \url{https://gretl.sourceforge.net}
		for details);
	\item A \cpp compiler with \texttt{-std=c++17} support (\texttt{g++},
		\texttt{clang++}, and \texttt{icpx} have been tested).
\end{enumerate}

Next, copy (or at least link to) all required plugins from the \gretl
installation directory (default location: \texttt{/usr/local/lib/gretl-gtk3/})
to \texttt{gretl/lib}.

Finally, to compile and install the module, run \texttt{meson setup builddir}
followed by \texttt{ninja -C builddir install} in the \gretlpy root directory.
This will generate the file \texttt{\_gretl.cpython-3xx-x86\_64-linux-gnu.so}
in the \texttt{builddir} directory, where \texttt{3.xx} refers to the version
of \python being used.

To verify the installation, run \texttt{python3 -c "import gretl; gretl.about(True)"};
this should produce output similar to Figure~\ref{fig:mswindows_about}.

